\newif\ifuseseminar
\useseminartrue
\input{../../latex_common/header.tex.frag}

\title{Mecânica Quântica II -- Partículas Idênticas e Transformações}

\begin{document}
\begin{enumerate}
	\item Considere um sistema de duas partículas {\bf idênticas} e responda às questões
	      abaixo:
	      \begin{enumerate}
		      \item Considere que o estado de uma partícula é descrito por $\ket{a} \in
			            \mathbb{V}$. Utilize o produto tensorial para definir os dois possíveis
		            espaços para os estados de duas partículas.
		      \item Definimos um estado separável de duas partículas como sendo da forma
		            $\ket{\psi} = \ket{a}\otimes\ket{b}$. Os estados que não podem ser
		            escritos dessa forma são chamados emaranhados. Os estados representando
		            duas partículas idênticas podem ser separáveis? Justifique sua resposta.
		      \item Com base nas respostas anteriores, explique como podemos realizar
		            experimentos em laboratório para medir o estado de uma partícula
		            $\ket{a}$ sem considerar todas as outras partículas idênticas presentes no
		            sistema. Para tanto, considere dois estados $\ket{T}$ e $\ket{L}$
		            representando pacotes gaussianos na Terra e na Lua, respectivamente, onde
		            $$\braket{x|T} =
			            \frac{1}{(2\pi)^{1/4}}\exp\left(-\frac{(x-\mu_T)^2}{4}\right),
			            \qquad \braket{x|L} =
			            \frac{1}{(2\pi)^{1/4}}\exp\left(-\frac{(x-\mu_L)^2}{4}\right),$$ com
		            $\mu_L\gg\mu_T$. Demonstre que tanto o estado simétrico quanto o
		            anti-simétrico de $\ket{L}$ e $\ket{T}$ são aproximadamente
		            separáveis quando calculamos a função de onda do conjunto em $(x_L,
			            x_T)$ onde $|x_T-\mu_L|\gg1$ e $|x_L-\mu_T|\gg1$.
		      \item Explique como podemos testar empiricamente se o estado de duas
		            partículas é simétrico ou anti-simétrico. Descreva como a simetria e a
		            anti-simetria estão relacionadas ao conceito de troca de partículas
		            idênticas e como as técnicas experimentais podem ser usadas para medir os
		            estados de duas partículas e verificar sua simetria.
	      \end{enumerate}
	\item Considere um sistema uni-dimensional com uma partícula e respondas às
	      questões abaixo:
	      \begin{enumerate}
		      \item Define o operador de translação $\hat{T}_\epsilon$ que desloca a
		            função de onda de uma partícula de uma distância $\epsilon$. Use
		            transformações ativas para definir o operador de translação.
		      \item Mostre que o gerador do operador de translação é o operador momento
		            linear $\hat{p}$.
		      \item Discuta a liberdade de fase associada ao operador de translação e
		            como ela está relacionada ao operador momento linear. Ou seja, se
		            definimos $\hat{T}_\epsilon\ket{x} = e^{i\epsilon
					            g(x)/\hbar}\ket{x+\epsilon}$, como podemos escolher $g(x)$ de forma
		            a garantir que $\hat{T}^\dagger_\epsilon\hat{p}\hat{T}_\epsilon =
			            \hat{p}$?
	      \end{enumerate}
	\item Em um sistema tridimensional com uma partícula, responda às questões sobre as
	      transformações discretas abaixo:
	      \begin{enumerate}
		      \item Defina o operador de reflexão $\hat{R}_x$ que inverte a coordenada
		            $x$. Use transformações ativas para definir o operador de reflexão.
		      \item O operador paridade $\hat{P}$ inverte todas as coordenadas espaciais.
		            Defina o operador paridade e mostre que ele é o produto de três operadores
		            de reflexão. Mostre também que o operador de reflexão combinado com uma
		            rotação de $\pi$ em torno de um eixo perpendicular ao plano de reflexão
		            é equivalente ao operador paridade.
		      \item Mostre que se a Hamiltoniana de um sistema é invariante sob rotações,
		            então invariância sob reflexões implica invariância sob paridade.
		      \item Encontre os autovalores e autovetores do operador de paridade para o
		            caso de uma partícula em uma dimensão. Mostre que os autovalores são
		            $\pm1$ e que os autovetores são funções pares e ímpares.
	      \end{enumerate}
	\item Inversão temporal é uma transformação que inverte o tempo. Considere um
	      sistema com uma partícula e responda às questões abaixo:
	      \begin{enumerate}
		      \item Defina o operador de inversão temporal $\hat{T}$ que inverte o tempo.
		            Usando a equação de Schrödinger para a função de onda de uma partícula,
		            mostre que a inversão temporal requer que a
		            função de onda seja transformada de acordo com
		            $\psi(x,t)\rightarrow\hat{T}\psi(x,t)=\psi^*(x,-t)$.
		      \item Usando o resultado acima, mostre que a inversão temporal transforma
		            a Hamiltoniana na representação de posição de acordo com
		            $\hat{T}^\dagger\hat{H}\hat{T} = \hat{H}^*$.
		      \item Mostre que a inversão temporal é uma transformação unitária e que
		            $\hat{T}^\dagger\hat{T} = \hat{I}$.
	      \end{enumerate}

\end{enumerate}

\end{document}
